\section{Foundations: theory}
\begin{frame}[plain]{Part I: Foundations}
\vspace{1em}\Large The canonical sharp RD design
\vspace{1.2em}\normalsize
\begin{enumerate}\setlength{\itemsep}{.65em}
\item The sharp RD design and its local effect
\item RD plots
\item Continuity based estimation and inference
\item Validation and falsification
\end{enumerate}
\F{Sections 1--5}
\end{frame}

\subsection{The sharp RD design}
\begin{frame}{The ingredients of a sharp RD design}
\begin{items}
\item Each unit has a score $X_i$ and faces a known cutoff $c$.
\item The assignment rule is $T_i=\ind\{X_i\geq c\}$.
\item In a \textbf{sharp} design, actual treatment equals assignment: $D_i=T_i$.
\item Foundations considers one continuous score, one cutoff, and perfect compliance.
\end{items}
\F{Section 2}
\end{frame}

\begin{frame}{Treatment assignment at the cutoff}
\begin{columns}
\begin{column}{.59\textwidth}\fig[.71]{sharp_assignment.pdf}\end{column}
\begin{column}{.37\textwidth}
The rule creates a discontinuous change in treatment probability.\\[1em]
Whether this change identifies a causal effect depends on assumptions about outcomes.
\end{column}
\end{columns}
\F{Section 2, Figure 1}
\end{frame}

\begin{frame}{Potential outcomes and observed outcomes}
For each unit, define $Y_i(1)$ and $Y_i(0)$.
\[
Y_i=T_iY_i(1)+(1-T_i)Y_i(0).
\]
\begin{items}
\item Only one potential outcome is observed for each unit.
\item Above the cutoff we observe treated outcomes. Below it we observe untreated outcomes.
\item Treatment effects require a comparison with the missing counterfactual.
\end{items}
\F{Section 2}
\end{frame}

\begin{frame}{The sharp RD effect}
Define the conditional mean potential outcomes
\[
\mu_d(x)=\E[Y_i(d)\mid X_i=x],\qquad d\in\{0,1\}.
\]
The target is the average treatment effect \textbf{at the cutoff}:
\[
\tau_{\mathrm{SRD}}=\E[Y_i(1)-Y_i(0)\mid X_i=c]
=\mu_1(c)-\mu_0(c).
\]
\begin{items}
\item For a continuous score, conditioning at $c$ describes a conditional mean function.
\item The target does not require a large group of observations exactly at $c$.
\end{items}
\F{Section 2}
\end{frame}

\begin{frame}{Continuity of the counterfactual mean functions}
\begin{columns}
\begin{column}{.57\textwidth}\fig[.7]{sharp_effect.pdf}\end{column}
\begin{column}{.40\textwidth}
Assume $\mu_0(x)$ and $\mu_1(x)$ are continuous at $c$.\\[.8em]
The solid portions are observable. The dashed portions are counterfactual.\\[.8em]
Continuity connects each observed limit to its missing value at the cutoff.
\end{column}
\end{columns}
\F{Section 2, Figure 2}
\end{frame}

\begin{frame}{Identification by one-sided limits}
On the two sides of the threshold,
\[
\E[Y_i\mid X_i=x]=
\begin{cases}\mu_0(x),&x<c,\\\mu_1(x),&x\geq c.\end{cases}
\]
Continuity gives
\[
\boxed{\tau_{\mathrm{SRD}}=
\lim_{x\downarrow c}\E[Y_i\mid X_i=x]
-\lim_{x\uparrow c}\E[Y_i\mid X_i=x].}
\]
\begin{items}
\item The observed outcome mean can jump because treatment changes.
\item Another intervention changing at the same cutoff can undermine this interpretation.
\end{items}
\F{Section 2}
\end{frame}

\begin{frame}{The local nature of the effect}
\begin{items}
\item The design identifies an effect for units at the assignment threshold.
\item Treatment effects may vary with the score and other characteristics.
\item A local effect can answer a relevant policy question about marginal eligibility.
\item Effects far from $c$, or under a substantially different cutoff, require additional assumptions.
\end{items}
\vspace{.5em}
\textbf{Identification and policy scope are separate questions.}
\F{Section 2.2}
\end{frame}

\begin{frame}{What continuity allows}
\begin{items}
\item The score can strongly predict both potential outcomes.
\item Treated and control units can differ substantially away from the threshold.
\item Continuity does not impose random assignment throughout a fixed interval.
\item In a finite sample, estimation must approximate the limits using nearby observations.
\end{items}
\F{Sections 2 and 4.1}
\end{frame}

\subsection{RD plots}
\begin{frame}{The purpose of an RD plot}
\begin{items}
\item Show the relationship between the score and the outcome on each side.
\item Reveal support, outliers, curvature, and a possible jump at $c$.
\item Combine \textbf{local averages in bins} with separate smooth fits.
\item Keep the graphical description distinct from formal estimation of the discontinuity.
\end{items}
\F{Section 3; \textcite{cct2015}}
\end{frame}

\begin{frame}{How binned means work}
Partition the score into bins $I_1,\ldots,I_J$. For bin $j$,
\[
\bar X_j=\frac{1}{n_j}\sum_{i:X_i\in I_j}X_i,
\qquad
\bar Y_j=\frac{1}{n_j}\sum_{i:X_i\in I_j}Y_i.
\]
\begin{items}
\item Plot one point $(\bar X_j,\bar Y_j)$ per bin.
\item Construct bins separately on each side so no bin mixes treatment statuses.
\item A point summarizes many units. It is not generally an individual observation.
\end{items}
\F{Section 3}
\end{frame}

\begin{frame}{Evenly spaced and quantile spaced bins}
\begin{columns}[T]
\begin{column}{.47\textwidth}
\textbf{Evenly spaced bins}\\[.5em]
Similar score width within each side.\\[.6em]
Numbers of observations can vary considerably.\\[.6em]
Sparse regions can produce noisy bin means.
\end{column}
\begin{column}{.47\textwidth}
\textbf{Quantile spaced bins}\\[.5em]
Similar numbers of observations per bin.\\[.6em]
Score widths vary with the density.\\[.6em]
Dense regions receive a finer score partition.
\end{column}
\end{columns}
\F{Section 3.1}
\end{frame}

\begin{frame}{The number of bins}
\begin{items}
\item Few wide bins smooth the plot heavily and can hide local variation.
\item Many narrow bins expose variation but make individual means less precise.
\item Data driven bin selection formalizes the tradeoff.
\item The two sides can use different numbers of bins.
\end{items}
\F{Section 3.2}
\end{frame}

\begin{frame}{IMSE and mimicking variance plots}
\begin{items}
\item \textbf{IMSE selection} balances integrated squared approximation error against variance when displaying the regression function.
\item \textbf{Mimicking variance selection} uses more bins to reveal the outcome's variability.
\item These are different graphical objectives and can produce different looking plots.
\item A smooth fitted curve can accompany either choice.
\end{items}
\F{Sections 3.2.1--3.2.2}
\end{frame}

\begin{frame}{Graphical choices and causal inference}
\begin{items}
\item Report the displayed score range and the binning rule.
\item Fit curves separately on either side of the cutoff.
\item A global polynomial used to describe a plot need not be the estimator used for causal inference.
\item Confidence intervals for bin means do not provide a confidence interval for the jump.
\end{items}
\F{Sections 3 and 4.1}
\end{frame}

\subsection{Continuity based estimation}
\begin{frame}{Local polynomial approximation}
\begin{columns}
\begin{column}{.58\textwidth}\fig[.70]{local_polynomial.pdf}\end{column}
\begin{column}{.39\textwidth}
Approximate each regression function in a neighborhood of $c$.\\[.8em]
Estimate its intercept at the cutoff.\\[.8em]
Subtract the left intercept from the right intercept.
\end{column}
\end{columns}
\F{Section 4.1, Figure 12}
\end{frame}

\begin{frame}{The local polynomial estimator}
Let $x_i=X_i-c$ and $r_p(x)=(1,x,\ldots,x^p)'$. Separately for $s\in\{-,+\}$,
\[
\widehat\beta_s=\arg\min_{\beta}\sum_{i\in\mathcal I_s}
 K\!\left(\frac{x_i}{h_s}\right)\big[Y_i-r_p(x_i)'\beta\big]^2,
\]
where $\mathcal I_- =\{i:x_i<0\}$ and $\mathcal I_+=\{i:x_i\geq0\}$.
\[
\widehat\tau=e_1'(\widehat\beta_+-\widehat\beta_-),
\qquad e_1=(1,0,\ldots,0)'.
\]
\F{Section 4.2}
\end{frame}

\begin{frame}{Local linear regression}
With $p=1$, the fitted lines are
\[
\widehat m_-(x)=\widehat\alpha_-+\widehat\beta_-x,
\qquad
\widehat m_+(x)=\widehat\alpha_++\widehat\beta_+x.
\]
\[
\widehat\tau=\widehat\alpha_+-\widehat\alpha_-.
\]
\begin{items}
\item The slopes can differ across sides.
\item Centering the score at $c$ makes each intercept a fitted cutoff value.
\item Local quadratic regression adds curvature without imposing a global fit.
\end{items}
\F{Section 4.2}
\end{frame}

\begin{frame}{A single regression representation}
Within a common window and using the same kernel weights,
\[
Y_i=\alpha+\tau T_i+\beta x_i+\gamma T_ix_i+\varepsilon_i.
\]
\begin{items}
\item Below the cutoff: intercept $\alpha$, slope $\beta$.
\item Above the cutoff: intercept $\alpha+\tau$, slope $\beta+\gamma$.
\item The coefficient on $T_i$ estimates the cutoff jump.
\item The algebraic representation alone does not provide bias corrected inference.
\end{items}
\F{Section 4.2}
\end{frame}

\begin{frame}{The kernel function}
Two common choices for the weights are
\[
K_U(u)=\ind\{|u|\leq1\},\qquad
K_T(u)=(1-|u|)\ind\{|u|\leq1\}.
\]
\begin{items}
\item A uniform kernel gives equal weight within the chosen window.
\item A triangular kernel gives more weight to observations near $c$.
\item Kernel weighting localizes the regression. It does not randomize treatment.
\end{items}
\F{Section 4.2.1}
\end{frame}

\begin{frame}{The polynomial order}
\begin{items}
\item Local linear regression is a practical default for estimating a level discontinuity.
\item Local quadratic regression can check sensitivity to curvature.
\item Higher orders can amplify variance and sensitivity near a boundary.
\item An elaborate polynomial over the full score range can give distant observations substantial influence on the cutoff fit.
\end{items}
\F{Section 4.2.1}
\end{frame}

\begin{frame}{Bandwidth and local sample size}
\[
\mathcal W_h=[c-h_-,c+h_+].
\]
\begin{items}
\item The bandwidth determines which observations inform the cutoff fit.
\item Smaller $h$ uses observations closer to $c$, reducing approximation error.
\item Larger $h$ generally improves precision by using more observations.
\item The effective sample sizes on each side matter more than the overall dataset size for precision at $c$.
\end{items}
\F{Section 4.2.2}
\end{frame}

\begin{frame}{Approximation bias near the boundary}
\fig[.70]{approximation_bias.pdf}
\src{A line can approximate a smooth curve differently as the estimation window changes.}
\F{Section 4.2.2, Figure 14}
\end{frame}

\begin{frame}{Bias, variance, and mean squared error}
For a level discontinuity, under regularity conditions and a common $h$,
\[
\operatorname{Bias}(\widehat\tau)\approx h^{p+1}\mathcal B,
\qquad
\V(\widehat\tau)\approx\frac{\mathcal V}{nh}.
\]
\[
\operatorname{MSE}(\widehat\tau)
\approx h^{2(p+1)}\mathcal B^2+\frac{\mathcal V}{nh}.
\]
\begin{items}
\item The squared bias and variance move in opposite directions.
\item For $p=1$, balancing their orders gives $h_{\mathrm{MSE}}\propto n^{-1/5}$.
\end{items}
\F{Sections 4.2.2--4.2.3; corrected bias order in the 2023 erratum}
\end{frame}

\begin{frame}{Data driven bandwidth selection}
\begin{items}
\item Feasible MSE selection estimates the unknown curvature and variance terms.
\item A common bandwidth and separate bandwidths are both possible.
\item The chosen bandwidth depends on the outcome, specification, and sample.
\item Selecting a bandwidth to obtain a desired sign or significance level invalidates the stated analysis.
\end{items}
\F{Sections 4.2.2--4.2.4}
\end{frame}

\subsection{Inference}
\begin{frame}{Why conventional inference can fail}
The conventional statistic treats approximation bias as negligible:
\[
T_{\mathrm{conv}}=\frac{\widehat\tau-\tau_0}{\widehat{\mathrm{se}}(\widehat\tau)}.
\]
\begin{items}
\item At an MSE optimal bandwidth, bias is generally of the same order as the standard error.
\item Centering a normal interval on $\widehat\tau$ can then produce poor coverage.
\item Heteroskedasticity robust standard errors address sampling variability, but do not remove smoothing bias.
\end{items}
\F{Section 4.3}
\end{frame}

\begin{frame}{Undersmoothing and bias correction}
\begin{items}
\item \textbf{Undersmoothing:} choose a smaller bandwidth so bias becomes negligible relative to sampling uncertainty.
\item This sacrifices precision and requires a defensible rule for how much to shrink $h$.
\item \textbf{Bias correction:} estimate the leading approximation bias and subtract it.
\item Estimating that bias introduces additional uncertainty.
\end{items}
\F{Section 4.3}
\end{frame}

\begin{frame}{Robust bias corrected inference}
Estimate the bias using a higher order local polynomial, often $q=p+1$:
\[
\widehat\tau_{\mathrm{bc}}=\widehat\tau-\widehat B.
\]
Use a variance estimate that includes the uncertainty in $\widehat B$:
\[
\mathrm{CI}_{\mathrm{RBC}}=
\left[\widehat\tau_{\mathrm{bc}}\ \pm\ z_{1-\alpha/2}\widehat{\mathrm{se}}_{\mathrm{RBC}}\right].
\]
\begin{items}
\item Correcting the center without correcting the standard error is insufficient.
\item The estimation bandwidth $h$ and bias bandwidth $b$ can differ.
\end{items}
\F{Section 4.3; \textcite{cct2014}}
\end{frame}

\begin{frame}{Point estimation and interval estimation}
\begin{items}
\item \textbf{MSE optimal bandwidths} target the accuracy of the point estimator.
\item \textbf{Coverage error optimal bandwidths} target the coverage accuracy of a confidence interval.
\item Robust bias correction allows valid asymptotic inference at an MSE optimal bandwidth under appropriate conditions.
\item A bandwidth optimized for one objective need not be optimal for the other.
\end{items}
\F{Section 4.3}
\end{frame}

\begin{frame}{Reading an RD results table}
\begin{items}
\item The conventional point estimate is the difference between the fitted cutoff intercepts.
\item The bias corrected estimate subtracts estimated smoothing bias.
\item The robust interval centers on the bias corrected estimate and uses the corrected variance.
\item A report should identify the estimator, interval method, polynomial order, kernel, bandwidths, and effective sample sizes.
\end{items}
\F{Section 4.3}
\end{frame}

\subsection{Additional practical issues}
\begin{frame}{Covariate adjustment}
\begin{items}
\item Predetermined covariates can improve precision when they predict the outcome.
\item A common additive adjustment can preserve the original RD effect under suitable balance and regularity conditions.
\item Treatment affected variables can introduce post-treatment bias.
\item Covariate adjustment cannot establish continuity when the underlying design is invalid.
\end{items}
\F{Section 4.4; \textcite{ccft2019}}
\end{frame}

\begin{frame}{Dependence and effective information}
\begin{items}
\item Standard errors must reflect the sampling and assignment environment.
\item Clustering can be appropriate when outcomes share shocks within sampling groups.
\item Few clusters near the cutoff can make conventional approximations unreliable.
\item Repeated score values introduce a separate support problem, developed in Extensions.
\end{items}
\F{Section 4.4}
\end{frame}

\subsection{Validation and falsification}
\begin{frame}{What validation can establish}
\begin{items}
\item The key continuity assumptions involve unobserved counterfactual outcomes.
\item Observable implications can expose reasons those assumptions are implausible.
\item Institutional knowledge is needed to interpret diagnostic evidence.
\item Passing a collection of tests supports credibility, but cannot prove identification.
\end{items}
\F{Section 5}
\end{frame}

\begin{frame}{Predetermined covariates and placebo outcomes}
\begin{items}
\item Use variables that the treatment could not have affected.
\item Estimate their discontinuities with the same care given to the main outcome.
\item Large or systematic jumps can indicate changes in composition or other confounders.
\item With many covariates, isolated small $p$-values can arise by chance. Their magnitude and substantive meaning matter.
\end{items}
\F{Section 5.1}
\end{frame}

\begin{frame}{The density of the running variable}
\begin{items}
\item A discontinuity in the score density can indicate precise sorting across the cutoff.
\item Modern density tests compare one-sided density estimates directly.
\item A density jump requires an explanation in the context of how the score is generated.
\item A smooth density alone cannot exclude every form of sorting or confounding.
\end{items}
\F{Section 5.2; \textcite{cj2020}}
\end{frame}

\begin{frame}{Placebo cutoffs}
\begin{items}
\item Estimate artificial discontinuities at points where the treatment rule does not change.
\item Restrict each analysis to observations on the same side of the actual cutoff.
\item A placebo window crossing the true treatment threshold can pick up a real effect.
\item Placebo jumps can reveal misspecification, other policies, or chance variation.
\end{items}
\F{Section 5.3}
\end{frame}

\begin{frame}{Observations closest to the cutoff}
\begin{items}
\item Heaping or manipulation may concentrate at the threshold itself.
\item A donut analysis removes observations very close to $c$.
\item It then requires more extrapolation to recover a cutoff effect.
\item Any change in results should be interpreted together with the reason for exclusion and the loss of local support.
\end{items}
\F{Section 5.4}
\end{frame}

\begin{frame}{Sensitivity to the bandwidth}
\begin{items}
\item Compare a planned range of smaller and larger bandwidths.
\item Examine magnitudes and confidence intervals, alongside sample sizes.
\item Report how the bias bandwidth changes when $h$ changes.
\item Significance across several correlated specifications is not a sequence of independent confirmations.
\end{items}
\F{Section 5.5}
\end{frame}

\begin{frame}{A complete sharp RD analysis}
\begin{enumerate}\setlength{\itemsep}{.55em}
\item Define the rule, outcome, and causal effect at the cutoff.
\item Inspect support and construct informative RD plots.
\item Estimate local polynomials with principled bandwidth selection.
\item Use inference that accounts for smoothing bias and dependence.
\item Assess predetermined outcomes, density, and specification sensitivity.
\item State the effect's local policy scope and remaining threats.
\end{enumerate}
\F{Sections 2--6}
\end{frame}
